\section{From decaying potentials to crossing comparisons}
\label{sec:crossings}

Propositions~\ref{prop:two-trajectories} and~\ref{prop:matching} give
potential arrays with small, decaying moment bounds on both lattices at the same physical
temperature $\beta_*$.  We now turn this information into three geometric
consequences: local conditional lower bounds, comparison of finite test
vectors on tori, and exclusion of interior four-passage bottlenecks.
The distinction between smallness and decay matters.  Smallness supplies
uniform passage estimates; decay makes the comparison errors vanish,
although $\sum_k X_k$ need not be finite.

The crossing propositions of~\cite[Section~7]{R} are stated along its
binary tuned trajectory.  Their proofs use the local scale construction
and independent primitive environment variables, together with
\begin{equation}
 \sup_k X_k\leq x_{\rm cr},\qquad X_k\longrightarrow0.
 \label{eq:crossing-hypothesis}
\end{equation}
We give the extension to the arrays used here, including the conditional
product estimate on which it depends.  Throughout, the arrays form an
actual trajectory of the exact local map of~\cite[Section~3]{R}, with
its conditional implementations and reversals specified below.  On a
torus we use the same local formulas and independent primitive variables,
so that every embedded stencil has its plane law.  Merely assuming the
displayed bounds for unrelated arrays would not suffice.
Throughout the section $s_k=s_0L^k$;
$\eta$ denotes the enlarged environment, including the independent
priorities and thresholds of the scale construction.  The Gibbs field
contains spins and a locally finite family of finite-valued auxiliary
variables, whose reference law is a product.  These auxiliary variables
are distinct from $\eta$; their decorated Gibbs law need not be a product.

\subsection{Buffered tests and the three passage estimates}

Draw the northwest--southeast diagonal in each lattice square and
interpolate spins affinely.  The zero contours have the prescribed
north-east/south-west pairing, up to a displacement of one mesh unit.
A colored path uses vertices with the specified sign and edges of this
triangulation; the diagonals introduce no additional Ising interactions.
A \emph{corridor} is a polygonal topological rectangle with two opposite
marked boundary arcs, called gates.  Its crossing test asks for a colored
path between the gates.  A \emph{band} is the closed region between two
disjoint simple polygons; its circuit test asks for a colored path
separating them.  We approximate these fixed nondegenerate polygons by
subcomplexes, preserving the order of marked sides, with an error of
order one lattice spacing.  Each test comes with a specified open
neighborhood in which spins, and the incident lattice edges, remain free.
This neighborhood is its buffer.

A passage between two terminal sets is a contour subarc joining them,
with interior in the prescribed region.  Distinct passages have disjoint
open parameter intervals on each contour, although endpoints at terminal
hits may coincide.  Thus several passages of one contour are counted
separately.  Cropping at the last hit of the first terminal before the
first hit of the second gives precisely this convention.

For $S_r(z,\alpha)=z+e^{i\alpha}[-r,r]^2$, consider the annulus
$A(z,\alpha;r,R)=S_R(z,\alpha)\setminus\operatorname{int}S_r(z,\alpha)$.
Write $F_4$ and $F_1$ for the events of at least four and at least one
passage across this annulus.  Write $H_2$ for two passages across it
confined to $z+e^{i\alpha}\{\operatorname{Im}w\geq0\}$.  The line in
this last definition is a geometric constraint inside a free bulk
chart; it carries no boundary pins.

We use the pure inputs of~\cite[Section~2, ``Iterated four-passage
bound'', ``Pure buffered RSW consequences'', and ``Two collars around
an annulus''; Section~5, ``Two bulk passages confined to a half-plane'']{S}.
For some $c_0>0$, they give
\begin{equation}
 \begin{aligned}
 \limsup_{\delta\downarrow0}\Prob_0(F_4(r,R))&\leq C(r/R)^{2+c_0},\\
 \limsup_{\delta\downarrow0}\Prob_0(F_1(r,R))&\leq C(r/R)^{c_0},\\
 \limsup_{\delta\downarrow0}\Prob_0(H_2(r,R))&=o(r/R).
 \end{aligned}
 \label{eq:crossing-pure-inputs}
\end{equation}
The ratio $r/R$ is fixed before the mesh limit.  Conditional laws with
arbitrary pins beyond proportional inner and outer free collars have
density at most a fixed factor $D$ relative to the pure bulk marginal;
$D$ does not grow with $R/r$.  These assertions permit the fixed crops,
rotations, and translations needed below.

Fix a gap factor $G$ that separates successive ring collars, including
their enlarged environmental neighborhoods.  Choose smaller collars
that will survive the local scale changes.  For a candidate ratio $a>1$
and each shape $i=4,h,1$,
let $q_i^0$ be a deterministic upper bound for its pure conditional
probability at ratio $a$, uniform over placements and arbitrary pins
beyond these smaller collars, once its inner lattice radius is
sufficiently large.  These bounds include the conditional density factor
and the fixed crop losses.  The pure estimates allow us to choose a
large fixed $a$, such bounds $q_i^0$, exponents, and tolerances satisfying
\begin{equation}
 p_4>2,\qquad p_h>1,\qquad p_1>0,\qquad
 q_i^0+\tau_i<(Ga)^{-p_i},\quad \tau_i>0,\quad i=4,h,1.
 \label{eq:crossing-slack}
\end{equation}
For $H_2$, first obtain strict slack below $(Ga)^{-1}$ from the little-oh
bound and then choose $p_h>1$ sufficiently close to one.  Require $s_0$
large enough that these pure bounds hold at finite lattice scales for
the three fixed shapes; this is the additional starting-scale choice
reserved in Section~\ref{sec:tuning}.  Slightly relaxed radii and compactness of directions and
normalized offsets give uniformity over their placements.  All choices
in this paragraph precede the choice of the disorder interval.

One deterministic input explains why exactly these three exponents
occur.  By~\cite[Section~5, ``Passages forced by a pivotal square'']{S},
if changing spins in a square of radius $b$ about $z$ changes a test $T$,
the common exterior contours force successive four-passage, half-plane
two-passage, and one-passage checks.  More precisely, let $a_T>0$ be
smaller than the relevant side lengths and separations, and put
\begin{equation}
 u=\min\{a_T,\max\{b,\dist(z,\mathcal E_T)\}\},\qquad
 v=\min\{a_T,\max\{b,\dist(z,\mathcal V_T)\}\},
 \label{eq:crossing-pivotal-distances}
\end{equation}
where $\mathcal E_T$ and $\mathcal V_T$ are its polygon sides and
vertices, including confinement walls for $H_2$.  The three checks
are forced between scales $b,u$, then $u,v$, then $v,a_T$.
A bounded number of rings may be omitted at transitions, and their
centers may move to the nearby side or vertex.  The implication permits
arbitrary fillings of the changed square.  For the three fixed annular
shapes, choose $a_T$ as a sufficiently small multiple of the inner
radius $r$, so every retained ring has inner radius at most $r/2$ and
its collars stay in the free buffer.

\begin{figure}[tb]
 \centering
 \begin{tikzpicture}[x=1cm,y=1cm,>=Stealth,
  frame/.style={draw=black!30,dashed,rounded corners=1pt},
  testedge/.style={draw=Orange!85!black,line width=1pt},
  passage/.style={draw=MidnightBlue,line width=.8pt},
  edit/.style={draw=Orange!85!black,fill=Orange!18,line width=.7pt}]
  \begin{scope}
    \node[font=\small\bfseries] at (1.75,2.2) {Bulk range};
    \draw[frame] (.55,.25) rectangle (2.95,1.85);
    \draw[edit] (1.55,.85) rectangle (1.95,1.25);
    \draw[passage] (1.55,.92) .. controls (1.12,.77) and (.91,.63) .. (.55,.53);
    \draw[passage] (1.59,1.25) .. controls (1.3,1.48) and (1.15,1.67) .. (1.02,1.85);
    \draw[passage] (1.92,1.25) .. controls (2.22,1.36) and (2.46,1.63) .. (2.63,1.85);
    \draw[passage] (1.95,.94) .. controls (2.28,.69) and (2.59,.69) .. (2.95,.57);
    \node[font=\scriptsize] at (1.75,1.04) {$z$};
    \node[font=\small,align=center] at (1.75,-.25) {$b<\ell<u$\\four passages};
  \end{scope}
  \begin{scope}[xshift=4.55cm]
    \node[font=\small\bfseries] at (1.75,2.2) {Side range};
    \draw[frame] (.55,.25) rectangle (2.95,1.85);
    \draw[testedge] (.30,.70)--(3.20,.70);
    \draw[edit] (1.58,.94) rectangle (1.86,1.22);
    \draw[passage] (1.45,.76) .. controls (1.02,.96) and (.82,1.42) .. (.72,1.85);
    \draw[passage] (1.97,.76) .. controls (2.30,1.01) and (2.60,1.39) .. (2.73,1.85);
    \draw[<->,thin] (2.80,.71)--(2.80,1.08);
    \node[font=\scriptsize,anchor=west] at (2.83,.90) {$u$};
    \node[font=\scriptsize,anchor=west,text=Orange!65!black] at (.32,.47) {test side};
    \node[font=\small,align=center] at (1.75,-.25) {$u<\ell<v$\\two confined passages};
  \end{scope}
  \begin{scope}[xshift=9.10cm]
    \node[font=\small\bfseries] at (1.75,2.2) {Vertex range};
    \draw[frame] (.55,.25) rectangle (2.95,1.85);
    \draw[testedge] (.35,.70)--(1.34,.70)--(1.34,1.97);
    \fill[Orange!85!black] (1.34,.70) circle (.04);
    \draw[edit] (1.67,.99) rectangle (1.95,1.27);
    \draw[passage] (1.96,1.13) .. controls (2.23,1.37) and (2.57,1.50) .. (2.95,1.67);
    \draw[<->,thin] (1.38,.68)--(1.81,1.07);
    \node[font=\scriptsize,anchor=north west] at (1.52,.82) {$v$};
    \node[font=\scriptsize,anchor=west,text=Orange!65!black] at (.35,.47) {test vertex};
    \node[font=\small,align=center] at (1.75,-.25) {$v<\ell<a_T$\\one passage};
  \end{scope}
  \draw[->,black!55] (1.10,-.98)--(12.30,-.98);
  \foreach \x/\lab in {1.15/b,4.18/u,8.45/v,12.25/a_T}{
    \draw[black!55] (\x,-.92)--(\x,-1.04);
    \node[below,font=\small] at (\x,-1.04) {$\lab$};
  }
  \node[font=\scriptsize,text=black!70] at (6.70,-1.59)
    {Increasing distance from the edit square};
\end{tikzpicture}

 \caption{The three geometric ranges in the pivotal-square implication.
 Orange squares indicate edits; orange sides and vertices belong to
 the test polygons, not to a pinned domain boundary.  Blue arcs represent
 the required exterior passages, schematically, at an intermediate
 scale $\ell$.  The actual checks use separated rings, may be recentered
 on a side or vertex, and omit rings near the transition scales.
 Dashed frames illustrate free neighborhoods, not literal ring placements.
 The distinction between bulk, side, and vertex locations yields the
 exponent requirements $p_4>2$, $p_h>1$, and $p_1>0$ in the location count.}
 \label{fig:pivotal-ranges}
\end{figure}
Figure~\ref{fig:pivotal-ranges} summarizes these ranges.  The scales
$u$ and $v$ locate the change from a two-dimensional set of possible
edit centers to neighborhoods of sides and then vertices.

\subsection{Local conditional laws and a single reversal}

Here and until physical lengths are explicitly introduced, distances
are in lattice units.  If $W$ is a free chart and $\zeta$ specifies the
exterior spin and reference auxiliary variables, the level-$k$
specification is
\begin{equation}
 \mu_{k,\eta,W}^{\zeta}(\dd\chi_W)
 =\frac{\exp F_{k,W}^{\zeta}(\chi_W)}
 {\mu_{{\rm pure},W}^{\zeta}(\exp F_{k,W}^{\zeta})}
 \mu_{{\rm pure},W}^{\zeta}(\dd\chi_W),
 \qquad
 F_{k,W}^{\zeta}=\sum_{i:\,\supp f_i^{(k)}\cap W\ne\varnothing}
 f_i^{(k)}(\chi_W\zeta_{W^c}).
 \label{eq:crossing-specification}
\end{equation}
The reference law is pure critical Ising times the product law of its
auxiliary variables.  At level zero, $\mu_{0,\eta,W}^{\zeta}$ is the
conditional nearest-neighbor Ising law under consideration.

We import the ``Local implementation'' and ``A coupling of consecutive
fields'' of~\cite[Section~3]{R}.  Their hypotheses are the array
hypotheses used in the scale construction: each potential is bounded,
real, spin-flip even and pure centered, with field support and primitive
environmental dependence within $Cs_k$ of its cell; the primitive
variables are independent; and the array is stationary on the scale
grid and square symmetric in law.  A finite first moment suffices for
the mark estimate below, and follows here from the fortieth moment.
No symmetry under negating a bond variable is required.

The local implementation has the following precise conditional scope.
With \emph{old} exterior field variables fixed, the step can be performed
inside $W$ and is exact.  Nontrivial fresh-coin decorations whose full
stencils leave $W$ are suppressed, their old logarithmic terms are
retained, and integrations whose full stencils leave $W$ are skipped.
All pairwise coin comparisons used to select separated parent cells
are retained, including comparisons that do not lead to an integration;
retained logarithmic terms are included in each permitted integration
they meet.  Every
nontrivially decorated fresh coin is therefore sampled inside $W$.
The output specification agrees with the full level-$(k+1)$
specification at distance at least $C_1s_{k+1}$ from $\partial W$.
Conditioning on the exterior of this eroded chart \emph{after} sampling
the output gives a mixture of its conditional specifications.  Exactness
does not assert an unchanged input kernel if the fresh coins are
independently pinned beforehand.

For fixed $\eta$, reversal changes spins only in marked squares of
radius $b_k=5s_{k+1}$ about parent centers.  The marks are independent
conditional on $\eta$ and independent of the entire output Gibbs
sample.  Choose the parameters themselves as exterior-independent
measurable functions of the local primitive-variable stencil, with
\begin{equation}
 p_{x,k}(\eta)\leq\min\left\{1,
 C_L\sum_{i:\,|s_ki-x|_\infty\leq C_Ls_{k+1}}\|f_i^{(k)}\|\right\},
 \qquad \E p_{x,k}\leq C_LX_k.
 \label{eq:crossing-marks}
\end{equation}
This applies also to the localized step and to unused candidate centers.
For example, if $S_x$ bounds the local logarithmic change, the old
conditional density is at least $e^{-S_x}$ times the new one.  Keeping
the new sample with probability $e^{-S_x}$ and otherwise sampling the
residual law proves the assertion, since $1-e^{-S_x}\leq\min\{1,S_x\}$.
Separation of the integration squares permits independent decisions.

For a check of inner radius $r$, start with
$W=S_{2ar}(z,\alpha)\setminus S_{r/2}(z,\alpha)$ and use
$W_k=W\ominus M_1s_k$, where erosion removes a layer along both
boundary components.  Choose
$M_1s_{k+1}\geq M_1s_k+C_1s_{k+1}$ and then $M_0$ large enough that
$r\geq M_0s_k$ leaves the smaller proportional collars fixed above.
Raise the first retained ring radius by a fixed factor if necessary so
all enlarged stencils fit in disjoint collars.

For marks $M$ from just one reversal and a spin sample $\sigma$, let
$\mathcal C(\sigma,M)$ be the configurations agreeing with $\sigma$
outside their marked squares.  Define
\[
 \begin{aligned}
 T^+(\sigma,M)&=\1\{T=1\text{ for some configuration in }
                         \mathcal C(\sigma,M)\},\\
 D_T(\sigma,M)&=\1\{T\text{ is nonconstant on }\mathcal C(\sigma,M)\}.
 \end{aligned}
\]
With $\pi_{k-1,\eta}$ the mark law for reversing step $k-1$, set
\begin{equation}
 B_{k,T}(\eta)=\sup_\zeta\mu_{k,\eta,W_k}^{\zeta}(T),\qquad
 B^+_{k,T}(\eta)=\sup_\zeta
       (\mu_{k,\eta,W_k}^{\zeta}\otimes\pi_{k-1,\eta})(T^+),\quad k\geq1.
 \label{eq:crossing-envelopes}
\end{equation}
Only $B_{0,T}$ is used at level zero.  Each supremum is a finite maximum:
only finitely many exterior variables can enter the chart, and their
alphabets are finite.  It is consequently measurable and depends only
on primitive environmental variables in the chart and a fixed
$O(s_k)$ neighborhood.  Taking the supremum \emph{before} environmental
expectation preserves this locality and bounds arbitrary mixtures of
exterior assignments: each fixed exterior kernel uses the same local
primitive variables, so maximizing over exterior values introduces no
additional environmental variables.

\begin{lemma}[Error from one reversal]
\label{lem:crossing-one-step}
Consider a fixed buffered test $T$ under a level-$k$ specification, with
marks from step $k-1$.  Suppose every retained ring check satisfies
$\E B^+_{k,T_i}\leq q_i^0+\tau_i$.  Retain rings only beyond a fixed
large multiple of $s_k$, with their collars enlarged by all field,
mark, and environmental stencils disjoint and avoiding the inserted
mark's primitive-variable stencil.  There is $c>0$, independent
of $T$, such that
\begin{equation}
 \E\sup_\zeta(\mu_{k,\eta}^{\zeta}\otimes\pi_{k-1,\eta})(D_T)
 \leq C_TX_{k-1}\min\{1,(s_k/a_T)^c\}.
 \label{eq:crossing-one-step}
\end{equation}
The same estimate bounds $T^+-T$, and holds for deterministic partial
activation of the marks.  Constants are uniform over the placements
and dilations of the three fixed check shapes.
\end{lemma}

\begin{proof}
Hold the output spin sample fixed and activate relevant squares in a
deterministic order.  The first square
whose addition permits two values of $T$ is pivotal for some assignment
inside the preceding squares.  Reverting only the new square in that
assignment shows that the passages forced by
\eqref{eq:crossing-pivotal-distances} are possible using previous marks.
Enlarging to all other marks in each ring can only increase this event.

For an inserted mark at $x$ and retained rings $T_1,\ldots,T_m$, condition
on all Gibbs variables outside one eroded ring chart and on all marks
whose squares do not meet its test support.  The enlarged collars put
every square relevant to this support inside the chart and make the
other ring events measurable under this conditioning.  The remaining
field has conditional law~\eqref{eq:crossing-specification}, with every potential
touching the chart included, and the remaining marks retain their
independent law.  This conditions on output auxiliary variables only
after their joint law has been constructed.  It does not condition on
the reconstructed old spins or on the other possibility events themselves.
Iterating gives
\begin{equation}
 \sup_\zeta\bigl[p_{x,k-1}(\eta)
       \Prob_{k,\eta}^{\zeta}(T_1^+\cap\cdots\cap T_m^+)\bigr]
 \leq p_{x,k-1}(\eta)\prod_{\ell=1}^m B^+_{k,T_\ell}(\eta).
 \label{eq:crossing-conditional-product}
\end{equation}
The ring neighborhoods are disjoint and avoid the primitive-variable
stencil of $p_{x,k-1}$.  Independence of these primitive variables,
including the priorities and thresholds, therefore gives
\begin{equation}
 \E\left[p_{x,k-1}\prod_{\ell=1}^m B^+_{k,T_\ell}\right]
 =\E p_{x,k-1}\prod_{\ell=1}^m\E B^+_{k,T_\ell}
 \leq C_LX_{k-1}\prod_{\ell=1}^m(q_{i(\ell)}^0+\tau_{i(\ell)}).
 \label{eq:crossing-independent-product}
\end{equation}
This factorization uses no property of the one-edge distribution beyond
primitive independence.  Unconditional quenched passage probabilities
would not provide the local factors needed for it.

Let $b=5s_k<a_T$.  Rings of ratio $a$ spaced in ratio $Ga$ and
\eqref{eq:crossing-slack} bound the last product, including the bounded
losses at transitions, by
\[
 C_T(b/u)^{p_4}(u/v)^{p_h}(v/a_T)^{p_1}.
\]
Normalize $a_T=1$.  Centers in a dyadic distance class $(u,v)$ occupy
area at most $C_Tuv$: near a vertex the relevant side length is $O_T(v)$
and its transverse width is $O(u)$; the class $u=v=1$ includes the
remaining interior.  Since centers have spacing comparable to $b$,
there are at most $C_Tuv/b^2$ such centers, including cell overhangs.
Their total contribution is at most
\[
 C_TX_{k-1}(b/u)^{p_4-2}(u/v)^{p_h-1}v^{p_1}.
\]
The three ratios multiply to $b$, so this is bounded by
$C_TX_{k-1}b^d$ for $d=\min\{p_4-2,p_h-1,p_1\}>0$.
There are $O_T((1+\log(1/b))^2)$ classes; absorbing this factor gives
$c=d/2$.  If $b\geq a_T$, only a bounded number of candidate squares
meet the fixed test, and the crude mark estimate suffices.  Finally,
$T^+-T\leq D_T$, and partial activation decreases the possibility events.
\end{proof}

\subsection{Closing the fixed-shape bootstrap}

\begin{lemma}[Uniform conditional passage bounds]
\label{lem:crossing-bootstrap}
There is a fixed $x_{\rm cr}>0$ such that $\sup_jX_j\leq x_{\rm cr}$
implies, for each of the three checks $T_i$ with inner radius
$r\geq M_0s_k$,
\begin{align}
 \E B_{k,T_i}
 &\leq q_i^0+C_{\rm b}\sum_{j=k}^{J(r)}X_j(s_j/r)^c,
 \label{eq:crossing-bootstrap-plain}\\
 \E B^+_{k,T_i}
 &\leq q_i^0+C_{\rm b}\sum_{j=k}^{J(r)}X_j(s_j/r)^c
               +C_{\rm b}X_{k-1}(s_k/r)^c,\qquad k\geq1,
 \label{eq:crossing-bootstrap-possible}
\end{align}
where $J(r)=\max\{j:M_0s_j\leq r\}$.  Both upper bounds are at most
$q_i^0+\tau_i$.  They hold for all placements, including embedded bulk
torus charts, with arbitrary exterior Gibbs data.
\end{lemma}

\begin{proof}
Use strong induction on dyadic ranges of $r/s_k$, simultaneously for
all levels and placements.  At terminal ratios
$M_0s_k\leq r<M_0s_{k+1}$ only a bounded number of potentials meet the
chart.  If $S$ is their total norm, the bounded logarithmic tilt in
\eqref{eq:crossing-specification} gives
\[
 \sup_\zeta\|\mu_{k,\eta,W_k}^{\zeta}
                  -\mu_{{\rm pure},W_k}^{\zeta}\|_{\TV}
 \leq\min\{1,2S\},\qquad \E S\leq C X_k.
\]
Since $(s_k/r)^c\geq(M_0L)^{-c}$, this proves the terminal plain bound
with a fixed $C_{\rm b}$.

For a nonterminal plain check, the local scale change expresses the new
interior law as a mixture of level-$(k+1)$ specifications.  Its plain
check has smaller ratio $r/s_{k+1}$.  Every ring controlling the reversal
has inner radius at most $r/2$, is evaluated at level $k+1$, and has
the required free collars.  Its ratio to its spacing lies in a strictly
smaller induction range.  The inductive possibility estimates therefore
justify Lemma~\ref{lem:crossing-one-step}, at cost
$CX_k(s_{k+1}/r)^c$.  This is a fixed multiple of the $j=k$ term in
\eqref{eq:crossing-bootstrap-plain}; the next-level plain bound supplies
the remaining terms.

After proving the plain bound at a given ratio, compare its possibility
event using marks from step $k-1$.  Its rings are now evaluated at level
$k$, still have inner radii at most $r/2$, and hence again belong to
strictly smaller induction ranges.  Lemma~\ref{lem:crossing-one-step}
adds $CX_{k-1}(s_k/r)^c$.  If no admissible ring fits, the test and edit
scales differ by a bounded factor, so a bounded count of candidate marks
gives the same estimate.  This also starts the terminal possibility
induction without a circular appeal to a passage estimate at its own
ratio.

The errors remain within the fixed tolerances because
\begin{equation}
 \sum_{j=k}^{J(r)}(s_j/r)^c
 \leq\frac{M_0^{-c}}{1-L^{-c}}.
 \label{eq:crossing-geometric-sum}
\end{equation}
Choose $C_{\rm b}$ to cover the terminal and one-step constants, and
then choose $x_{\rm cr}$ so that $C_{\rm b}x_{\rm cr}$ times this bound,
plus the additional single-level term, is below $\min_i\tau_i$.
These choices use only the three fixed check shapes and the fixed scale
parameters.  Each invocation of the induction uses precisely these
tolerances, which closes both estimates.
\end{proof}

Lemma~\ref{lem:crossing-one-step} now applies to every prescribed finite
family of buffered polygonal tests.  Such a family changes its constants
$C_T,a_T$, but does not change $x_{\rm cr}$.  The remaining convergence
arguments use the elementary consequence of geometric spacing
\begin{equation}
 \sum_{j<K}X_j\min\{1,(s_{j+1}/r)^c\}+X_K\longrightarrow0
 \label{eq:crossing-weighted-decay}
\end{equation}
as $r,K\to\infty$, provided $s_K/r$ stays between fixed positive bounds.
Indeed the sum of the weights is uniformly bounded.  For $j\geq m$
bound $X_j$ by $\sup_{j\geq m}X_j$, and for the finitely many $j<m$
let $r\to\infty$.  Then let $m\to\infty$.  This is where decay, rather
than summability, is sufficient.

\subsection{The three geometric consequences}

\begin{proposition}[Uniform local conditional comparison]
\label{prop:local-comparison}
Assume the array hypotheses above and~\eqref{eq:crossing-hypothesis}.
Fix a bounded polygonal window with a finite family $\mathcal T$ of
corridor crossings, band circuits, or passage checks, strictly inside
specified free buffers.  Annular windows may have an inner and an outer
free collar.  Dilate the whole geometry by $r\to\infty$ in lattice units.
There is an exterior-independent measurable envelope $\mathcal E_r(\eta)$
with $\E\mathcal E_r\to0$, uniformly over translations and scale-grid
positions, such that every level-zero conditional test-vector law is
within $\mathcal E_r(\eta)$ in total variation of a mixture of pure
conditional test-vector laws on a still-buffered interior window.

Consequently, if an event in the test vector has pure conditional
probability at least $c_*>0$ under all exterior assignments after a
sufficiently small fixed erosion of the buffers, its level-zero
probability is at least any $c<c_*$, uniformly over exterior assignments,
outside an environmental event of probability tending to zero.  For a
nearest-neighbor initial law, the event that this worst-boundary lower
bound holds depends only on the bonds in the window and its incident
layer.
\end{proposition}

\begin{proof}
Choose a terminal level with $s_K/r$ between two sufficiently small
fixed positive constants.  The cumulative erosion has width $O(s_K)$,
so all tests retain their prescribed buffers.  At each reversal,
\eqref{eq:crossing-conditional-product} supplies an error bound for
every exterior assignment before taking any mixture.  Concretely it is
the sum over inserted centers of their mark parameters times the
products of the appropriate ring envelopes, or just the mark parameters
when no ring fits.  Summing over the finite family and steps gives an
envelope whose expectation is bounded by
\begin{equation}
 \E\mathcal E_r\leq C_{\mathcal T}
 \left(\sum_{j<K}X_j\min\{1,(s_{j+1}/r)^c\}+X_K\right).
 \label{eq:crossing-local-error}
\end{equation}
The final term comes from deleting the bounded number of level-$K$
potentials touching the remaining window, using the logarithmic-tilt
bound in the preceding proof.  Conditional on the final exterior this
leaves the pure law, and disintegration gives the stated mixture.
Equation~\eqref{eq:crossing-weighted-decay} proves convergence.

Every member of the pure mixture gives the selected event probability
at least $c_*$.  Thus every initial exterior assignment gives probability
at least $c_*-\mathcal E_r$.  Average this inequality over the algorithmic
choices conditional on the original bond environment $\omega$, then
apply Markov's inequality.  The resulting exceptional probability is
at most $\E\mathcal E_r/(c_*-c)$.  Although the bounding envelope may use
more variables, the actual worst-boundary probability in the initial
nearest-neighbor law uses only the window bonds and incident layer;
its lower-bound event is therefore local.  Uniformity means a bound for
each placement with the same error rate, not simultaneous success at
all placements.
\end{proof}

\begin{proposition}[Expected total variation for torus tests]
\label{prop:torus-tests}
Assume the same array hypotheses and~\eqref{eq:crossing-hypothesis}.
Fix finitely many buffered tests $\mathcal T$ in a common planar chart,
now in physical units.  Let $\delta\downarrow0$ and choose tori of lattice
period $Ms_K$, with $M$ a sufficiently large fixed integer, $K\to\infty$,
and physical periods $\delta Ms_K$ in a fixed compact subinterval of
$(0,\infty)$ large enough for the chart and the complete scale stencils
to embed.  Give the torus iid bonds with the original law.  If
$\nu_{\omega}^{\mathcal T}$ and $\nu_{\rm pure}^{\mathcal T}$ are the
initial quenched and pure critical laws of the whole Boolean test vector,
then
\begin{equation}
 \E_{\omega}\|\nu_{\omega}^{\mathcal T}
                     -\nu_{\rm pure}^{\mathcal T}\|_{\TV}\longrightarrow0.
 \label{eq:crossing-torus-tv}
\end{equation}
Such tori can be chosen for every deterministic mesh sequence.
\end{proposition}

\begin{proof}
Embedded stencils have exactly their plane primitive-variable laws and
the disjointness properties used above.  Couple consecutive levels at
fixed enlarged environment.  The expected error for test $T$ at step
$j$ is at most
\[
 C_TX_j\min\{1,(\delta s_{j+1}/a_T^{\rm phys})^c\},
\]
by Lemmas~\ref{lem:crossing-one-step} and~\ref{lem:crossing-bootstrap}.
Their sum tends to zero by the same geometric-weight argument as
\eqref{eq:crossing-weighted-decay}.  At level $K$ the torus contains
$O(M^2)$ potentials, whose expected total norm is $O_M(X_K)$.
Deleting them couples the entire terminal law to the pure torus with
vanishing expected total variation.  Glue these couplings and take a
union bound over steps and the fixed test list.  No estimate on the union
of marks from different levels is needed.

The initial test law depends only on $\omega$, whereas the final pure
law is deterministic.  Averaging the coupling bounds over algorithmic
choices therefore proves exactly~\eqref{eq:crossing-torus-tv}, rather
than merely comparison of annealed test laws.  Finally, choosing $K$ as
the first scale with $s_K\geq A\delta^{-1}$, for a sufficiently large
fixed $A$, places the physical periods between $MA$ and $MAL$.
\end{proof}

\begin{corollary}[No interior four-passage bottlenecks]
\label{cor:four-passages}
At $\beta_*$, let $F$ be a fixed compact set with positive clearance
inside the free part of the approximating domains, and choose a
sufficiently small fixed physical radius $R>0$ so that the required
annular charts about $F$ remain free.  Uniformly over all exterior spin
assignments,
\begin{equation}
 \lim_{r\downarrow0}\limsup_{\delta\downarrow0}
 \E_{\omega}\mu_{\omega}\bigl(
 \text{four contour passages from }r\text{ to }R
                       \text{ about some }z\in F\bigr)=0.
 \label{eq:crossing-four-passages}
\end{equation}
Passages have the parameter-interval multiplicities specified above.
Fixed-factor changes of radii and local contour redrawing errors are
permitted after corresponding crops.
\end{corollary}

\begin{proof}
At a fixed center, pack rings of ratio $a$, separated by factor $G$,
between a fixed multiple of $r$ and a fixed fraction of $R$.  For fixed
physical $r>0$, their lattice radii exceed $M_0s_0$ once the mesh is
small.  Conditioning off their charts as in
\eqref{eq:crossing-conditional-product}, and then applying
\eqref{eq:crossing-independent-product} without an inserted mark,
bounds the passage probability by the product of the plain envelopes.
Their expectations are at most $q_4^0+\tau_4<(Ga)^{-p_4}$.
Consequently the mesh upper limit is at most $C(r/R)^{p_4}$, uniformly
over the ambient exterior data.

Cover $F$ by an $r$-grid with $O_F(r^{-2})$ centers.  Enlarging the inner
radius and decreasing the outer radius by fixed factors makes four
passages about any point imply a grid check.  The union bound gives
$C_{F,R}r^{p_4-2}$ after the mesh limit.  Since $p_4>2$, this tends to
zero.  The same crops absorb lattice rounding and the local matched
drawings of the resolved contours.
\end{proof}

The initializations and their scale outputs satisfy every array
hypothesis used here.  Proposition~\ref{prop:two-trajectories} makes
the uniform bound on $X_k$ in~\eqref{eq:crossing-hypothesis} as small as needed
by reducing $\varepsilon_0(\rho)$, and makes $X_k$ decay.  The fixed
threshold $x_{\rm cr}$ depends only on the three check shapes and the
scale parameters.  Thus this reduction of $\varepsilon_0(\rho)$ precedes
the choice of a final domain, test family, or approximation sequence.
Proposition~\ref{prop:matching} places these conclusions on both lattices
at $\beta_*$.  The local comparison will give FK annular barriers; the
torus comparison and four-passage exclusion will control the interface.
